% Ad Atticum 9.5 --- headnote
% ad-atticum-09-05, 10 March 49 BC, Formian villa

\textit{Cicero to Atticus, written from the Formian villa on
10 March 49 BC (the manuscript dateline: \textit{Scr.\ im
Formiano vi Id.\ Mart.\ a.\ 705 (49)}). The letter is
formally earlier than the one that, in the OCT, precedes
it (9.4, of 12 March): the editorial order is canonical,
not chronological. The occasion is a letter from Atticus
sent on his own birthday, full of careful counsel about
where Cicero should go --- Adriatic or Tyrrhenian, Arpinum
or Formiae, each route with its own scandal.}

\textit{Section 1 reports the day's company at the villa:
Postumus and Quintus Fufius, both hurrying south to
Brundisium and both denouncing Pompey and the Senate ---
``if I cannot bear these things in my own villa, am I to
bear Curtius in the Senate-house?'' Section 2 turns to
the underlying question. Cicero declares the
commonwealth lost. He has been angry with Pompey, more
angry than with Caesar himself, because the causes of
events move him more than the events: he counts back ten
years of Pompeian negligence, the year of his own exile
foremost, and likens that day's anger to the Roman
memory of the disaster on the Allia, more mourned than
the Gallic sack itself. Section 3 then turns. Pompey's
benefits and his \textit{dignitas} are remembered;
Balbus's reports have made it clear that Caesar's whole
campaign aims at Pompey's destruction; and so Cicero
casts himself as Achilles answering Thetis after the
death of Patroclus --- the celebrated Homer quotation
(Iliad 18.96, 98--99), turned into the language of
\textit{officium}: such duties are worth purchasing with
life itself. Section 4 closes on the optimates: he has
no confidence in them, watches them give themselves over
to Caesar, and waits for news from Brundisium.}
